% !TEX program = xelatex
\documentclass[bwprint]{report}
\usepackage{tikz}
\usetikzlibrary{arrows.meta,positioning,shapes.geometric}

\title{软件系统综合课程设计练习2单元测试报告}
\stunum{待填写}
\stuname{程泽凯}
\stuclass{软件工程专业 2023 级}
\supervisor{课程指导组}
\dateinput{2026年6月17日}
\def\makeothertitle{}
\setlength{\cftsecnumwidth}{3.2em}
\setlength{\cftsubsecnumwidth}{4.4em}
\setlength{\cftsubsubsecnumwidth}{5.6em}
\newcolumntype{P}[1]{>{\raggedright\arraybackslash}p{#1}}
\lstset{basicstyle=\ttfamily\small,breaklines=true,frame=single,columns=fullflexible,keepspaces=true}
\sloppy

\begin{document}
\maketitle
\newpage
\pagestyle{plain}
\pagenumbering{Roman}
\begin{abstract}
这次练习 2 我主要做的是服务层代码的单元测试设计。被测项目仍然是“西南交通大学校园论坛系统”，测试对象选取了本人负责的 3 个代码单元。报告中保留核心代码、控制流图、基本路径、条件-判定覆盖、条件组合覆盖和 JUnit 5 示例，同时补充火车票退票业务的黑盒场景法用例。

\vspace{0.5em}
\noindent{\heiti 关键词：}单元测试；白盒测试；基本路径法；条件组合；JUnit；退票场景法
\end{abstract}
\tableofcontents
\newpage
\setcounter{page}{1}
\pagenumbering{arabic}

\section{实验过程简述}
我先从项目后端代码中挑出 3 个比较适合做白盒测试的单元，再把每个单元里的判断条件和提前返回分支标出来。随后根据控制流图拆基本路径，用 JUnit 5 和 Mockito 设计对应的测试方法。这个过程里没有把页面点击当成单元测试结果，而是尽量回到服务层方法本身，检查异常类型、返回对象、仓储调用和字段变化。

\section{个人测试点总表}
{
\setlength{\tabcolsep}{2pt}
\begin{longtable}{P{0.8cm}P{2.7cm}P{5.4cm}P{5.1cm}}
\toprule
序号 & 代码单元 & 方法位置 & 测试重点与覆盖目标\\
\midrule
\endfirsthead
\toprule
序号 & 代码单元 & 方法位置 & 测试重点与覆盖目标\\
\midrule
\endhead
1 & 后台封禁用户逻辑（后台用户管理） & Admin\hspace{0pt}Ops\hspace{0pt}Service. ban\hspace{0pt}User (Long,  String) & 重点：最后管理员保护和封禁日志；覆盖目标：保护分支和正常封禁分支都覆盖；方法：基本路径法、条件-判定覆盖、条件组合覆盖 \\
2 & 帖子举报创建逻辑（内容举报） & Report\hspace{0pt}Service. create\hspace{0pt}Post\hspace{0pt}Report (Long,  Create\hspace{0pt}Report\hspace{0pt}Request,  String) & 重点：原因合法性、自举报、重复举报；覆盖目标：治理入口的拒绝分支和保存分支都覆盖；方法：基本路径法、条件-判定覆盖、条件组合覆盖 \\
3 & 通知已读标记逻辑（通知中心） & Notification\hspace{0pt}Service. mark\hspace{0pt}Read (Long,  String) & 重点：通知存在性、归属校验、已读保存；覆盖目标：不泄露他人通知并能标记本人通知；方法：基本路径法、条件-判定覆盖、条件组合覆盖 \\
\bottomrule
\end{longtable}
}
这 3 个单元的共同点是：代码里都有明确的分支判断，而且每个判断都会影响业务结果。把它们拿来做练习，比只测试简单 getter 或纯数据对象更有意义。

\section{白盒测试设计与分析}
本节按照“核心代码段、控制流图、基本路径法、条件-判定覆盖、条件组合覆盖、JUnit 示例”的顺序展开。每个单元内部结构保持一致，但测试说明按各自业务分支单独写。

\subsection{后台封禁用户逻辑}
\textbf{被测代码单元：}\texttt{AdminOpsService.banUser(Long, String)}。该单元包含保留最后一个可用管理员的保护条件，适合做条件组合覆盖。

这段代码属于“后台用户管理”功能，测试时我主要盯住“最后管理员保护和封禁日志”。一开始看代码行数不算多，但异常分支会提前结束流程，如果只测成功路径，很容易把真正容易出错的地方漏掉。

\textbf{核心代码段摘录：}
\begin{lstlisting}[language=Java]
public User banUser(Long userId, String operatorName) {
    User user = userRepository.findById(userId)
            .orElseThrow(() -> new NotFoundException("user not found"));
    if ("ADMIN".equals(user.getRole()) && Boolean.TRUE.equals(user.getEnabled())
            && userRepository.countByRoleAndEnabledTrue("ADMIN") <= 1) {
        throw new ApiException(400, "keep at least one admin");
    }
    user.setEnabled(false);
    User saved = userRepository.save(user);
    User operator = userService.getUserByUsername(operatorName);
    log(operator, "BAN_USER", "USER", userId, "ban user");
    return saved;
}
\end{lstlisting}

\textbf{控制流图说明：}控制流图里把普通处理语句合并成顺序节点，把会改变执行方向的 if、orElseThrow、短路返回作为判定节点。这样画出来不追求每一行代码一个节点，但能看清测试时必须覆盖的分支。

\begin{figure}[H]
\centering
\resizebox{0.78\textwidth}{!}{%
\begin{tikzpicture}[node distance=6mm and 11mm, >=Stealth]
\tikzstyle{startstop}=[rectangle, rounded corners, minimum width=2.35cm, minimum height=0.68cm, text centered, draw=black, fill=gray!12, font=\normalsize]
\tikzstyle{process}=[rectangle, minimum width=2.8cm, minimum height=0.68cm, text centered, text width=3.15cm, draw=black, fill=blue!6, font=\normalsize]
\tikzstyle{decision}=[diamond, aspect=2.0, text centered, text width=2.8cm, draw=black, fill=orange!10, font=\normalsize]
\tikzstyle{arrow}=[->, thick]
\node (start) [startstop] {Start};
\node (d1) [decision, below=of start] {用户是否存在?};
\draw [arrow] (start) -- (d1);
\node (b1) [process, right=of d1] {用户不存在};
\draw [arrow] (d1) -- node[above,font=\small] {否} (b1);
\node (d2) [decision, below=of d1] {是否唯一可用管理员?};
\node (b2) [process, right=of d2] {拒绝封禁};
\draw [arrow] (d2) -- node[above,font=\small] {是} (b2);
\node (final) [process, below=of d2] {封禁用户并记录日志};
\draw [arrow] (d2) -- node[left,font=\small] {否} (final);
\node (end) [startstop, below=of final] {End};
\draw [arrow] (final) -- (end);
\draw [arrow] (d1) -- node[left,font=\small] {是} (d2);
\draw [arrow] (b1) |- (end);
\draw [arrow] (b2) |- (end);
\end{tikzpicture}%
}
\caption{控制流图}
\end{figure}

\textbf{基本路径法分析：}用户存在和最后管理员保护条件中三个子条件，共规划 5 条路径。 本单元的路径设计不是把所有输入都穷举一遍，而是把会让程序提前结束、进入异常处理或完成保存的路线分开。保护分支和正常封禁分支都覆盖。

{
\setlength{\tabcolsep}{3pt}
\begin{longtable}{P{1.2cm}P{12.0cm}}
\toprule
路径 & 业务含义\\
\midrule
\endfirsthead
\toprule
路径 & 业务含义\\
\midrule
\endhead
P1 & 用户不存在。 \\
P2 & 目标是唯一可用管理员，拒绝封禁。 \\
P3 & 目标是普通用户，封禁成功。 \\
P4 & 目标是管理员但仍有其他可用管理员，封禁成功。 \\
P5 & 目标管理员已禁用，再次封禁保持禁用并记录操作。 \\
\bottomrule
\end{longtable}
}
这些路径合在一起，能够覆盖主流程和主要异常出口。后续如果代码增加新的 if 分支，需要先回到这里补路径，再补测试方法。

\textbf{基本路径法测试用例：}下表把每条路径落成可执行的 JUnit 用例，重点是让 mock 条件和预期分支一一对应。
{
\setlength{\tabcolsep}{3pt}
\begin{longtable}{P{2.4cm}P{4.9cm}P{5.8cm}}
\toprule
编号 & 测试数据与条件 & 预期结果\\
\midrule
\endfirsthead
\toprule
编号 & 测试数据与条件 & 预期结果\\
\midrule
\endhead
BAN-\hspace{0pt}P1 & user\hspace{0pt}Id 不存在 & Not\hspace{0pt}Found\hspace{0pt}Exception \\
BAN-\hspace{0pt}P2 & role=\hspace{0pt}ADMIN, enabled=\hspace{0pt}true, count=\hspace{0pt}1 & Api\hspace{0pt}Exception 400 \\
BAN-\hspace{0pt}P3 & role=\hspace{0pt}USER & enabled=\hspace{0pt}false，写操作日志 \\
BAN-\hspace{0pt}P4 & role=\hspace{0pt}ADMIN, enabled=\hspace{0pt}true, count=\hspace{0pt}2 & enabled=\hspace{0pt}false \\
BAN-\hspace{0pt}P5 & role=\hspace{0pt}ADMIN, enabled=\hspace{0pt}false & 不触发最后管理员保护 \\
\bottomrule
\end{longtable}
}
这些用例主要用来卡住执行路径。实际写代码时，最容易出问题的是 mock 条件没有配全，导致测试没有走到预想分支。

\textbf{条件-判定覆盖：}本部分关注每个判定条件至少取到一次真值和一次假值。对于短路判断，报告中会把未执行条件单独说明，不把它误写成已经覆盖。
{
\setlength{\tabcolsep}{3pt}
\begin{longtable}{P{2.4cm}P{5.0cm}P{5.8cm}}
\toprule
编号 & 覆盖条件 & 判定结果\\
\midrule
\endfirsthead
\toprule
编号 & 覆盖条件 & 判定结果\\
\midrule
\endhead
BAN-\hspace{0pt}CD1 & is\hspace{0pt}Admin=\hspace{0pt}T & 管理员条件取真 \\
BAN-\hspace{0pt}CD2 & enabled=\hspace{0pt}T & 启用状态条件取真 \\
BAN-\hspace{0pt}CD3 & last\hspace{0pt}Admin=\hspace{0pt}T & 保护条件取真 \\
BAN-\hspace{0pt}CD4 & 组合整体为假 & 允许封禁 \\
\bottomrule
\end{longtable}
}
条件-判定覆盖的作用是提醒我不要只看返回值，也要看每个 if 的真假两面有没有被触发。

\textbf{条件组合覆盖：}理论组合不一定都可执行，例如前一个判定已经抛异常时，后面的条件不会再执行。因此这里优先列出业务上可行、能被 JUnit 稳定构造的组合。
{
\setlength{\tabcolsep}{3pt}
\begin{longtable}{P{2.4cm}P{2.0cm}P{2.0cm}P{2.0cm}P{4.1cm}}
\toprule
编号 & 条件1 & 条件2 & 条件3/4 & 组合期望\\
\midrule
\endfirsthead
\toprule
编号 & 条件1 & 条件2 & 条件3/4 & 组合期望\\
\midrule
\endhead
BAN-\hspace{0pt}C1 & T & T & T & 拒绝封禁 \\
BAN-\hspace{0pt}C2 & T & T & F & 允许封禁 \\
BAN-\hspace{0pt}C3 & T & F & - & 允许封禁 \\
BAN-\hspace{0pt}C4 & F & - & - & 允许封禁 \\
\bottomrule
\end{longtable}
}
组合覆盖比单独的条件覆盖更容易暴露短路逻辑。例如前置条件为真时，后续仓储查询或密码判断可能根本不会发生，测试中应使用 \texttt{verify(..., never())} 或异常断言来确认。

\textbf{JUnit 5 测试类示例：}下面是按本单元整理的测试类骨架。导入语句按项目实际包名补齐即可，个别实体构造方法也可以换成项目已有的测试工厂方法。
\begin{lstlisting}[language=Java]
@ExtendWith(MockitoExtension.class)
class AdminBanUserWhiteBoxTest {
    @Mock UserRepository userRepository;
    @Mock OperationLogRepository operationLogRepository;
    @Mock UserService userService;
    @InjectMocks AdminOpsService adminOpsService;

    // 根据实际项目类名补齐 DTO 和实体的构造细节。
    private User user(String name) {
        User user = new User();
        user.setId(1L);
        user.setUsername(name);
        user.setEnabled(true);
        return user;
    }

@Test
void banUserKeepsLastEnabledAdmin() {
    admin.setRole("ADMIN");
    admin.setEnabled(true);
    when(userRepository.countByRoleAndEnabledTrue("ADMIN")).thenReturn(1L);
    assertThrows(ApiException.class, () -> adminOpsService.banUser(admin.getId(), "root"));
}

@Test
void banUserDisablesNormalUserAndLogsOperation() {
    normal.setRole("USER");
    User saved = adminOpsService.banUser(normal.getId(), "root");
    assertThat(saved.getEnabled()).isFalse();
    verify(operationLogRepository).save(any(OperationLog.class));
}
}
\end{lstlisting}

\textbf{缺陷与回归关注：}未发现缺陷。该保护条件很关键，后续修改用户角色逻辑时应优先回归。 这部分没有发现阻塞缺陷，但后续回归时还要注意 DTO 校验、异常提示和前端展示是否与服务层状态保持一致。


\subsection{帖子举报创建逻辑}
\textbf{被测代码单元：}\texttt{ReportService.createPostReport(Long, CreateReportRequest, String)}。该单元防止非法原因、自举报和重复举报，属于后台治理流程入口。

这段代码属于“内容举报”功能，测试时我主要盯住“原因合法性、自举报、重复举报”。一开始看代码行数不算多，但异常分支会提前结束流程，如果只测成功路径，很容易把真正容易出错的地方漏掉。

\textbf{核心代码段摘录：}
\begin{lstlisting}[language=Java]
public Report createPostReport(Long postId, CreateReportRequest request, String username) {
    User reporter = userService.getUserByUsername(username);
    Post post = postRepository.findById(postId)
            .orElseThrow(() -> new NotFoundException("post not found"));
    String reason = request.getReason().trim();
    if (!ALLOWED_REASONS.contains(reason)) {
        throw new ApiException(400, "bad reason");
    }
    if (post.getAuthor() != null && post.getAuthor().getId().equals(reporter.getId())) {
        throw new ApiException(400, "cannot report yourself");
    }
    if (reportRepository.existsByReporterAndPostAndStatus(reporter, post, STATUS_PENDING)) {
        throw new ApiException(400, "duplicate report");
    }
    Report report = new Report();
    report.setReporter(reporter);
    report.setPost(post);
    report.setReason(reason);
    return reportRepository.save(report);
}
\end{lstlisting}

\textbf{控制流图说明：}控制流图里把普通处理语句合并成顺序节点，把会改变执行方向的 if、orElseThrow、短路返回作为判定节点。这样画出来不追求每一行代码一个节点，但能看清测试时必须覆盖的分支。

\begin{figure}[H]
\centering
\resizebox{0.78\textwidth}{!}{%
\begin{tikzpicture}[node distance=6mm and 11mm, >=Stealth]
\tikzstyle{startstop}=[rectangle, rounded corners, minimum width=2.35cm, minimum height=0.68cm, text centered, draw=black, fill=gray!12, font=\normalsize]
\tikzstyle{process}=[rectangle, minimum width=2.8cm, minimum height=0.68cm, text centered, text width=3.15cm, draw=black, fill=blue!6, font=\normalsize]
\tikzstyle{decision}=[diamond, aspect=2.0, text centered, text width=2.8cm, draw=black, fill=orange!10, font=\normalsize]
\tikzstyle{arrow}=[->, thick]
\node (start) [startstop] {Start};
\node (d1) [decision, below=of start] {帖子是否存在?};
\draw [arrow] (start) -- (d1);
\node (b1) [process, right=of d1] {帖子不存在};
\draw [arrow] (d1) -- node[above,font=\small] {否} (b1);
\node (d2) [decision, below=of d1] {举报原因是否合法?};
\node (b2) [process, right=of d2] {拒绝非法原因};
\draw [arrow] (d2) -- node[above,font=\small] {否} (b2);
\node (d3) [decision, below=of d2] {是否举报自己的帖子?};
\node (b3) [process, right=of d3] {拒绝自举报};
\draw [arrow] (d3) -- node[above,font=\small] {是} (b3);
\node (d4) [decision, below=of d3] {是否已有待处理举报?};
\node (b4) [process, right=of d4] {拒绝重复举报};
\draw [arrow] (d4) -- node[above,font=\small] {是} (b4);
\node (final) [process, below=of d4] {保存举报记录};
\draw [arrow] (d4) -- node[left,font=\small] {否} (final);
\node (end) [startstop, below=of final] {End};
\draw [arrow] (final) -- (end);
\draw [arrow] (d1) -- node[left,font=\small] {是} (d2);
\draw [arrow] (b1) |- (end);
\draw [arrow] (d2) -- node[left,font=\small] {是} (d3);
\draw [arrow] (b2) |- (end);
\draw [arrow] (d3) -- node[left,font=\small] {否} (d4);
\draw [arrow] (b3) |- (end);
\draw [arrow] (b4) |- (end);
\end{tikzpicture}%
}
\caption{控制流图}
\end{figure}

\textbf{基本路径法分析：}帖子存在、原因合法、是否举报自己、是否重复四个判定，V(G)=5。 本单元的路径设计不是把所有输入都穷举一遍，而是把会让程序提前结束、进入异常处理或完成保存的路线分开。治理入口的拒绝分支和保存分支都覆盖。

{
\setlength{\tabcolsep}{3pt}
\begin{longtable}{P{1.2cm}P{12.0cm}}
\toprule
路径 & 业务含义\\
\midrule
\endfirsthead
\toprule
路径 & 业务含义\\
\midrule
\endhead
P1 & 帖子不存在。 \\
P2 & 举报原因不在允许集合内。 \\
P3 & 举报自己的帖子。 \\
P4 & 同一用户对同一帖子已有待处理举报。 \\
P5 & 合法举报保存成功。 \\
\bottomrule
\end{longtable}
}
这些路径合在一起，能够覆盖主流程和主要异常出口。后续如果代码增加新的 if 分支，需要先回到这里补路径，再补测试方法。

\textbf{基本路径法测试用例：}下表把每条路径落成可执行的 JUnit 用例，重点是让 mock 条件和预期分支一一对应。
{
\setlength{\tabcolsep}{3pt}
\begin{longtable}{P{2.4cm}P{4.9cm}P{5.8cm}}
\toprule
编号 & 测试数据与条件 & 预期结果\\
\midrule
\endfirsthead
\toprule
编号 & 测试数据与条件 & 预期结果\\
\midrule
\endhead
REPORT-\hspace{0pt}P1 & post\hspace{0pt}Id 不存在 & Not\hspace{0pt}Found\hspace{0pt}Exception \\
REPORT-\hspace{0pt}P2 & reason=随便写 & Api\hspace{0pt}Exception 400 \\
REPORT-\hspace{0pt}P3 & reporter=\hspace{0pt}post.author & Api\hspace{0pt}Exception 400 \\
REPORT-\hspace{0pt}P4 & exists pending=\hspace{0pt}true & Api\hspace{0pt}Exception 400 \\
REPORT-\hspace{0pt}P5 & reason=垃圾广告，无重复 & 保存 Report \\
\bottomrule
\end{longtable}
}
这些用例主要用来卡住执行路径。实际写代码时，最容易出问题的是 mock 条件没有配全，导致测试没有走到预想分支。

\textbf{条件-判定覆盖：}本部分关注每个判定条件至少取到一次真值和一次假值。对于短路判断，报告中会把未执行条件单独说明，不把它误写成已经覆盖。
{
\setlength{\tabcolsep}{3pt}
\begin{longtable}{P{2.4cm}P{5.0cm}P{5.8cm}}
\toprule
编号 & 覆盖条件 & 判定结果\\
\midrule
\endfirsthead
\toprule
编号 & 覆盖条件 & 判定结果\\
\midrule
\endhead
REPORT-\hspace{0pt}CD1 & post\hspace{0pt}Exists=\hspace{0pt}F & 帖子不存在 \\
REPORT-\hspace{0pt}CD2 & reason\hspace{0pt}Allowed=\hspace{0pt}F & 原因不合法 \\
REPORT-\hspace{0pt}CD3 & self\hspace{0pt}Report=\hspace{0pt}T & 拒绝自举报 \\
REPORT-\hspace{0pt}CD4 & duplicate=\hspace{0pt}T & 拒绝重复举报 \\
REPORT-\hspace{0pt}CD5 & 全部满足 & 创建成功 \\
\bottomrule
\end{longtable}
}
条件-判定覆盖的作用是提醒我不要只看返回值，也要看每个 if 的真假两面有没有被触发。

\textbf{条件组合覆盖：}理论组合不一定都可执行，例如前一个判定已经抛异常时，后面的条件不会再执行。因此这里优先列出业务上可行、能被 JUnit 稳定构造的组合。
{
\setlength{\tabcolsep}{3pt}
\begin{longtable}{P{2.4cm}P{2.0cm}P{2.0cm}P{2.0cm}P{4.1cm}}
\toprule
编号 & 条件1 & 条件2 & 条件3/4 & 组合期望\\
\midrule
\endfirsthead
\toprule
编号 & 条件1 & 条件2 & 条件3/4 & 组合期望\\
\midrule
\endhead
REPORT-\hspace{0pt}C1 & F & - & - & -；帖子不存在 \\
REPORT-\hspace{0pt}C2 & T & F & - & -；非法原因 \\
REPORT-\hspace{0pt}C3 & T & T & T & -；自举报 \\
REPORT-\hspace{0pt}C4 & T & T & F & T；重复举报 \\
REPORT-\hspace{0pt}C5 & T & T & F & F；成功 \\
\bottomrule
\end{longtable}
}
组合覆盖比单独的条件覆盖更容易暴露短路逻辑。例如前置条件为真时，后续仓储查询或密码判断可能根本不会发生，测试中应使用 \texttt{verify(..., never())} 或异常断言来确认。

\textbf{JUnit 5 测试类示例：}下面是按本单元整理的测试类骨架。导入语句按项目实际包名补齐即可，个别实体构造方法也可以换成项目已有的测试工厂方法。
\begin{lstlisting}[language=Java]
@ExtendWith(MockitoExtension.class)
class ReportCreateWhiteBoxTest {
    @Mock ReportRepository reportRepository;
    @Mock PostRepository postRepository;
    @Mock OperationLogRepository operationLogRepository;
    @Mock ReportReviewLogRepository reportReviewLogRepository;
    @Mock UserService userService;
    @Mock NotificationService notificationService;
    @InjectMocks ReportService reportService;

    // 根据实际项目类名补齐 DTO 和实体的构造细节。
    private User user(String name) {
        User user = new User();
        user.setId(1L);
        user.setUsername(name);
        user.setEnabled(true);
        return user;
    }

@Test
void createReportRejectsSelfReport() {
    post.setAuthor(reporter);
    request.setReason("垃圾广告");
    assertThrows(ApiException.class,
            () -> reportService.createPostReport(1L, request, "alice"));
}

@Test
void createReportSavesValidReport() {
    request.setReason("虚假信息");
    Report saved = reportService.createPostReport(1L, request, "bob");
    assertThat(saved.getReason()).isEqualTo("虚假信息");
    verify(reportRepository).save(any(Report.class));
}
}
\end{lstlisting}

\textbf{缺陷与回归关注：}未发现缺陷。建议前端原因下拉框和后端允许集合保持同步，减少非法原因提交。 这部分没有发现阻塞缺陷，但后续回归时还要注意 DTO 校验、异常提示和前端展示是否与服务层状态保持一致。


\subsection{通知已读标记逻辑}
\textbf{被测代码单元：}\texttt{NotificationService.markRead(Long, String)}。该单元用 NotFoundException 隐藏他人通知存在性，适合验证权限分支。

这段代码属于“通知中心”功能，测试时我主要盯住“通知存在性、归属校验、已读保存”。一开始看代码行数不算多，但异常分支会提前结束流程，如果只测成功路径，很容易把真正容易出错的地方漏掉。

\textbf{核心代码段摘录：}
\begin{lstlisting}[language=Java]
public Notification markRead(Long id, String username) {
    User recipient = requireUser(username);
    Notification notification = notificationRepository.findById(id)
            .orElseThrow(() -> new NotFoundException("notification not found"));
    if (notification.getRecipient() == null
            || !notification.getRecipient().getId().equals(recipient.getId())) {
        throw new NotFoundException("notification not found");
    }
    notification.setIsRead(true);
    return notificationRepository.save(notification);
}
\end{lstlisting}

\textbf{控制流图说明：}控制流图里把普通处理语句合并成顺序节点，把会改变执行方向的 if、orElseThrow、短路返回作为判定节点。这样画出来不追求每一行代码一个节点，但能看清测试时必须覆盖的分支。

\begin{figure}[H]
\centering
\resizebox{0.78\textwidth}{!}{%
\begin{tikzpicture}[node distance=6mm and 11mm, >=Stealth]
\tikzstyle{startstop}=[rectangle, rounded corners, minimum width=2.35cm, minimum height=0.68cm, text centered, draw=black, fill=gray!12, font=\normalsize]
\tikzstyle{process}=[rectangle, minimum width=2.8cm, minimum height=0.68cm, text centered, text width=3.15cm, draw=black, fill=blue!6, font=\normalsize]
\tikzstyle{decision}=[diamond, aspect=2.0, text centered, text width=2.8cm, draw=black, fill=orange!10, font=\normalsize]
\tikzstyle{arrow}=[->, thick]
\node (start) [startstop] {Start};
\node (d1) [decision, below=of start] {用户是否存在?};
\draw [arrow] (start) -- (d1);
\node (b1) [process, right=of d1] {用户不存在};
\draw [arrow] (d1) -- node[above,font=\small] {否} (b1);
\node (d2) [decision, below=of d1] {通知是否存在?};
\node (b2) [process, right=of d2] {通知不存在};
\draw [arrow] (d2) -- node[above,font=\small] {否} (b2);
\node (d3) [decision, below=of d2] {通知是否属于当前用户?};
\node (b3) [process, right=of d3] {按通知不存在处理};
\draw [arrow] (d3) -- node[above,font=\small] {否} (b3);
\node (final) [process, below=of d3] {设置已读并保存};
\draw [arrow] (d3) -- node[left,font=\small] {是} (final);
\node (end) [startstop, below=of final] {End};
\draw [arrow] (final) -- (end);
\draw [arrow] (d1) -- node[left,font=\small] {是} (d2);
\draw [arrow] (b1) |- (end);
\draw [arrow] (d2) -- node[left,font=\small] {是} (d3);
\draw [arrow] (b2) |- (end);
\draw [arrow] (b3) |- (end);
\end{tikzpicture}%
}
\caption{控制流图}
\end{figure}

\textbf{基本路径法分析：}用户存在、通知存在、通知归属三个判定，V(G)=4。 本单元的路径设计不是把所有输入都穷举一遍，而是把会让程序提前结束、进入异常处理或完成保存的路线分开。不泄露他人通知并能标记本人通知。

{
\setlength{\tabcolsep}{3pt}
\begin{longtable}{P{1.2cm}P{12.0cm}}
\toprule
路径 & 业务含义\\
\midrule
\endfirsthead
\toprule
路径 & 业务含义\\
\midrule
\endhead
P1 & 用户名不存在。 \\
P2 & 通知不存在。 \\
P3 & 通知存在但不属于当前用户。 \\
P4 & 通知属于当前用户，标记已读成功。 \\
\bottomrule
\end{longtable}
}
这些路径合在一起，能够覆盖主流程和主要异常出口。后续如果代码增加新的 if 分支，需要先回到这里补路径，再补测试方法。

\textbf{基本路径法测试用例：}下表把每条路径落成可执行的 JUnit 用例，重点是让 mock 条件和预期分支一一对应。
{
\setlength{\tabcolsep}{3pt}
\begin{longtable}{P{2.4cm}P{4.9cm}P{5.8cm}}
\toprule
编号 & 测试数据与条件 & 预期结果\\
\midrule
\endfirsthead
\toprule
编号 & 测试数据与条件 & 预期结果\\
\midrule
\endhead
NOTICE-\hspace{0pt}P1 & username 不存在 & Not\hspace{0pt}Found\hspace{0pt}Exception \\
NOTICE-\hspace{0pt}P2 & notification\hspace{0pt}Id 不存在 & Not\hspace{0pt}Found\hspace{0pt}Exception \\
NOTICE-\hspace{0pt}P3 & recipient\hspace{0pt}Id != current\hspace{0pt}User\hspace{0pt}Id & Not\hspace{0pt}Found\hspace{0pt}Exception \\
NOTICE-\hspace{0pt}P4 & recipient\hspace{0pt}Id == current\hspace{0pt}User\hspace{0pt}Id & is\hspace{0pt}Read=\hspace{0pt}true 并保存 \\
\bottomrule
\end{longtable}
}
这些用例主要用来卡住执行路径。实际写代码时，最容易出问题的是 mock 条件没有配全，导致测试没有走到预想分支。

\textbf{条件-判定覆盖：}本部分关注每个判定条件至少取到一次真值和一次假值。对于短路判断，报告中会把未执行条件单独说明，不把它误写成已经覆盖。
{
\setlength{\tabcolsep}{3pt}
\begin{longtable}{P{2.4cm}P{5.0cm}P{5.8cm}}
\toprule
编号 & 覆盖条件 & 判定结果\\
\midrule
\endfirsthead
\toprule
编号 & 覆盖条件 & 判定结果\\
\midrule
\endhead
NOTICE-\hspace{0pt}CD1 & user\hspace{0pt}Exists=\hspace{0pt}F & 用户不存在 \\
NOTICE-\hspace{0pt}CD2 & notice\hspace{0pt}Exists=\hspace{0pt}F & 通知不存在 \\
NOTICE-\hspace{0pt}CD3 & owner\hspace{0pt}Match=\hspace{0pt}F & 不泄露他人通知 \\
NOTICE-\hspace{0pt}CD4 & owner\hspace{0pt}Match=\hspace{0pt}T & 标记已读 \\
\bottomrule
\end{longtable}
}
条件-判定覆盖的作用是提醒我不要只看返回值，也要看每个 if 的真假两面有没有被触发。

\textbf{条件组合覆盖：}理论组合不一定都可执行，例如前一个判定已经抛异常时，后面的条件不会再执行。因此这里优先列出业务上可行、能被 JUnit 稳定构造的组合。
{
\setlength{\tabcolsep}{3pt}
\begin{longtable}{P{2.4cm}P{2.0cm}P{2.0cm}P{2.0cm}P{4.1cm}}
\toprule
编号 & 条件1 & 条件2 & 条件3/4 & 组合期望\\
\midrule
\endfirsthead
\toprule
编号 & 条件1 & 条件2 & 条件3/4 & 组合期望\\
\midrule
\endhead
NOTICE-\hspace{0pt}C1 & F & - & - & 用户不存在 \\
NOTICE-\hspace{0pt}C2 & T & F & - & 通知不存在 \\
NOTICE-\hspace{0pt}C3 & T & T & F & 他人通知，仍返回不存在 \\
NOTICE-\hspace{0pt}C4 & T & T & T & 标记成功 \\
\bottomrule
\end{longtable}
}
组合覆盖比单独的条件覆盖更容易暴露短路逻辑。例如前置条件为真时，后续仓储查询或密码判断可能根本不会发生，测试中应使用 \texttt{verify(..., never())} 或异常断言来确认。

\textbf{JUnit 5 测试类示例：}下面是按本单元整理的测试类骨架。导入语句按项目实际包名补齐即可，个别实体构造方法也可以换成项目已有的测试工厂方法。
\begin{lstlisting}[language=Java]
@ExtendWith(MockitoExtension.class)
class NotificationMarkReadWhiteBoxTest {
    @Mock NotificationRepository notificationRepository;
    @Mock UserRepository userRepository;
    @InjectMocks NotificationService notificationService;

    // 根据实际项目类名补齐 DTO 和实体的构造细节。
    private User user(String name) {
        User user = new User();
        user.setId(1L);
        user.setUsername(name);
        user.setEnabled(true);
        return user;
    }

@Test
void markReadDoesNotExposeOtherUsersNotification() {
    notification.setRecipient(otherUser);
    assertThrows(NotFoundException.class,
            () -> notificationService.markRead(9L, "alice"));
}

@Test
void markReadSetsReadFlag() {
    notification.setRecipient(alice);
    Notification saved = notificationService.markRead(9L, "alice");
    assertThat(saved.getIsRead()).isTrue();
    verify(notificationRepository).save(notification);
}
}
\end{lstlisting}

\textbf{缺陷与回归关注：}未发现缺陷。该分支体现了权限隐藏策略，需要和接口状态码契约保持一致。 这部分没有发现阻塞缺陷，但后续回归时还要注意 DTO 校验、异常提示和前端展示是否与服务层状态保持一致。


\section{测试环境与执行记录}
\begin{itemize}
\item 操作系统：Windows。
\item Java 环境：JDK 17，Spring Boot 3.2.0。
\item 单元测试框架：JUnit 5 + Mockito + Maven Surefire。
\item 实际执行命令：进入 \texttt{jiaodaluntan/backend} 后运行 \texttt{mvn test}。
\item 本次执行结果：通过 85，失败 0，错误 0，跳过 0。
\item 覆盖率工具：本轮未启用 JaCoCo，因此报告只记录 Surefire 的真实单元测试结果，不编造覆盖率百分比。
\end{itemize}

{
\setlength{\tabcolsep}{3pt}
\begin{longtable}{P{2.8cm}P{10.0cm}}
\toprule
项目 & 记录\\
\midrule
\endfirsthead
\toprule
项目 & 记录\\
\midrule
\endhead
执行命令 & mvn test \\
执行目录 & jiaodaluntan/\hspace{0pt}backend \\
测试源码位置 & src/\hspace{0pt}test/\hspace{0pt}java/\hspace{0pt}com/\hspace{0pt}swjtu/\hspace{0pt}forum \\
本报告对应测试类 & Admin\hspace{0pt}Ops\hspace{0pt}Service\hspace{0pt}Test（封禁用户）、Report\hspace{0pt}Service\hspace{0pt}Test（举报创建）、Notification\hspace{0pt}Service\hspace{0pt}Test（已读通知） \\
测试类数量 & 18 个 \\
测试方法数量 & 85 个 \\
成功/失败/错误/跳过 & 通过 85，失败 0，错误 0，跳过 0 \\
执行时间 & 12.139 s \\
完成时间 & 2026-\hspace{0pt}06-\hspace{0pt}17 16:01:22 +08:00 \\
Surefire 报告 & target/\hspace{0pt}surefire-\hspace{0pt}reports，已生成 18 组测试报告 \\
Ja\hspace{0pt}Co\hspace{0pt}Co 报告 & 本轮未启用 Ja\hspace{0pt}Co\hspace{0pt}Co；未编造覆盖率百分比 \\
备注 & 主要验证服务层分支逻辑，页面截图不作为本次单元测试证据 \\
\bottomrule
\end{longtable}
}

\section{缺陷记录与回归建议}
{
\setlength{\tabcolsep}{3pt}
\begin{longtable}{P{2.1cm}P{2.7cm}P{7.6cm}}
\toprule
类别 & 位置 & 说明\\
\midrule
\endfirsthead
\toprule
类别 & 位置 & 说明\\
\midrule
\endhead
已发现问题 & 本次 3 个服务层单元 & 本轮 mvn test 共 85 个测试全部通过，未记录到阻塞级服务层缺陷 \\
潜在风险 & 输入 DTO 校验 & 空用户名、非法邮箱、弱密码、空评论等最好在 DTO 或校验器层单独补测试，避免服务层假设输入永远合法 \\
潜在风险 & 登录限流/时间窗口 & 限流逻辑依赖时间或缓存状态，建议把 Login\hspace{0pt}Attempt\hspace{0pt}Service 单独写测试，不要只在登录方法里间接覆盖 \\
潜在风险 & 分页边界 & 分页工具类建议补 Integer.MAX\_VALUE、Integer.MIN\_VALUE 等极端值 \\
潜在风险 & 火车票退票 & 退票费取整规则容易写错，特别是 2.5 角和 7.5 角附近 \\
回归建议 & 测试目录 & 本轮测试已落到真实 src/\hspace{0pt}test/\hspace{0pt}java 目录；后续修改服务层后应继续执行 mvn test \\
\bottomrule
\end{longtable}
}
这部分没有为了凑数量去写虚假的 Bug。能确定的问题和风险分开记录，后续回归时更容易判断哪些需要马上修，哪些只是需要加测试保护。


\section{黑盒测试补充：火车票退票场景法}
\subsection{临近发车与办理渠道切换}
我的黑盒测试假设票价为 214.50 元，主线放在临近发车时的渠道切换，同时检查大额票价下手续费进位规则。车票开车时间固定为 2026 年 6 月 30 日 07:52。除特别说明外，场景默认张三未打印纸质车票；为了测试手续费取整规则，最后三条额外设置了不同票价。

\textbf{基本流：}张三进入退票流程，系统读取开车时间、退票申请时间、票价和纸质票状态；先判断办理渠道，再判断手续费比例，最后按 5 角规则取整并给出退款金额。

\textbf{备选流：}如果退票时间正好落在 15 天、48 小时、24 小时或 2 小时边界，需要按需求文字解释边界；如果已打印纸质车票，即使时间满足网上退票，也只能到车站窗口办理。

\textbf{异常流：}不足开车前 2 小时仍从网上入口提交时，系统应提示不能网上办理；手续费取整遇到 2.5 角和 7.5 角附近时，不能简单四舍五入到整数。

\begin{lstlisting}
S0 start
S1 read ticket: departure=2026-06-30 07:52, price, paperTicket
S2 if before 2 hours and no paper ticket -> online; else station window
S3 if >15 days -> 0%; else if >=48h -> 5%; else if >=24h -> 10%; else -> 20%
S4 round fee by 0.5 yuan rule
S5 refund = price - rounded fee
S6 output channel, fee, refund amount
\end{lstlisting}

\textbf{场景矩阵：}本矩阵先列业务分支，再展开成测试用例，避免只按时间随便挑几条。
{
\setlength{\tabcolsep}{3pt}
\begin{longtable}{P{2.3cm}P{4.6cm}P{4.7cm}}
\toprule
场景类型 & 覆盖点 & 说明\\
\midrule
\endfirsthead
\toprule
场景类型 & 覆盖点 & 说明\\
\midrule
\endhead
时间边界 & 15天、48小时、24小时、2小时 & 这些点最容易因为大于/大于等于写错而出错 \\
办理渠道 & 未打印且不少于2小时、已打印、不足2小时 & 渠道判断和手续费比例不是同一个判断 \\
费用取整 & 小于2.5角、2.5到7.5角、7.5角以上 & 需要额外票价才能把取整边界测出来 \\
\bottomrule
\end{longtable}
}

{
\setlength{\tabcolsep}{3pt}
\begin{longtable}{P{1.1cm}P{3.05cm}P{2.25cm}P{3.45cm}P{3.05cm}}
\toprule
编号 & 时间与票价 & 纸质票与渠道 & 场景说明 & 预期结果\\
\midrule
\endfirsthead
\toprule
编号 & 时间与票价 & 纸质票与渠道 & 场景说明 & 预期结果\\
\midrule
\endhead
S01 & 2026-\hspace{0pt}06-\hspace{0pt}14 07:51；票价 214.50 & 未打印；网上可退 & 大于15天 & 手续费0.00，退款214.50 \\
S02 & 2026-\hspace{0pt}06-\hspace{0pt}15 07:52；票价 214.50 & 未打印；网上可退 & 正好15天，按48小时以上处理 & 手续费10.50，退款204.00 \\
S03 & 2026-\hspace{0pt}06-\hspace{0pt}20 09:00；票价 214.50 & 未打印；网上可退 & 48小时以上且不超过15天 & 手续费10.50，退款204.00 \\
S04 & 2026-\hspace{0pt}06-\hspace{0pt}28 07:52；票价 214.50 & 未打印；网上可退 & 正好48小时 & 手续费10.50，退款204.00 \\
S05 & 2026-\hspace{0pt}06-\hspace{0pt}28 08:10；票价 214.50 & 未打印；网上可退 & 不足48小时且超过24小时 & 手续费21.50，退款193.00 \\
S06 & 2026-\hspace{0pt}06-\hspace{0pt}29 07:52；票价 214.50 & 未打印；网上可退 & 正好24小时，按24小时以上处理 & 手续费21.50，退款193.00 \\
S07 & 2026-\hspace{0pt}06-\hspace{0pt}29 22:00；票价 214.50 & 未打印；网上可退 & 不足24小时但超过2小时 & 手续费43.00，退款171.50 \\
S08 & 2026-\hspace{0pt}06-\hspace{0pt}30 05:52；票价 214.50 & 未打印；网上可退 & 正好2小时 & 手续费43.00，退款171.50 \\
S09 & 2026-\hspace{0pt}06-\hspace{0pt}30 06:10；票价 214.50 & 未打印；窗口办理 & 不足2小时 & 手续费43.00，退款171.50 \\
S10 & 2026-\hspace{0pt}06-\hspace{0pt}20 10:00；票价 214.50 & 已打印；窗口办理 & 时间满足网上退票但已取票 & 手续费10.50，退款204.00 \\
S11 & 2026-\hspace{0pt}06-\hspace{0pt}20 10:30；票价 101.00 & 未打印；网上可退 & 补充票价：尾数小于2.5角 & 5\%为5.05，取整5.00，退款96.00 \\
S12 & 2026-\hspace{0pt}06-\hspace{0pt}20 11:00；票价 105.00 & 未打印；网上可退 & 补充票价：尾数2.5角到7.5角 & 5\%为5.25，取整5.50，退款99.50 \\
S13 & 2026-\hspace{0pt}06-\hspace{0pt}20 11:30；票价 115.50 & 未打印；网上可退 & 补充票价：尾数7.5角以上 & 5\%为5.775，取整6.00，退款109.50 \\
\bottomrule
\end{longtable}
}

这组用例按本人对退票业务的理解独立设计。前十条主要卡时间和渠道边界，后三条专门卡手续费尾数取整，避免把金额规则只测成“按比例乘一下”。


\section{实验收获}
后台类代码最需要测试的是保护条件。封禁用户如果没有最后管理员保护，系统可能把自己锁死；举报如果不防重复和自举报，管理队列会很乱；通知如果把他人的通知返回成“无权访问”，反而可能泄露通知存在。练习路径法时，我觉得这些分支比成功路径更有价值，因为它们是系统保持可管理和可恢复的底线。

这次练习让我更明显地感觉到，白盒测试不是看页面能不能点，而是要把代码里的判断条件拆出来。控制流图虽然画起来有点麻烦，但它能逼着我检查“是不是只测了成功分支”。像登录限流、重复签到、库存不足这类路径，如果不画出来，很容易在写用例时被忽略。

JUnit 自动化对回归也很有帮助。手工看一遍代码只能说明当时理解了逻辑，测试方法保留下来以后，后面改服务层、改 DTO 或改仓储接口时，还能重新跑一遍确认分支没有被改坏。黑盒退票作业则提醒我，业务规则不一定难在流程长，而是难在边界：正好 48 小时、正好 2 小时、手续费尾数取整，这些点比普通样例更容易暴露问题。

\end{document}
